\section{Large Language Models}
\label{sec:LLMs}

Here I provide a brief overview of large language models (LLMs), their architecture, training methods, and capabilities. This is not meant to be an exhaustive review, but rather a primer to understand the context of LLM alignment.

\subsection{What are Language Models?}

Language models (LMs) are a class of machine learning models that are designed to understand and generate human language. They are typically probabilistic models that learn to predict the next token in a sequence of text given the preceding context. By training on large datasets of text data (corpora), language models can learn the statistical patterns and structures inherent in human language, enabling them to generate coherent and contextually relevant text.

Language models see the world as \term{tokens} which are generated by taking a string and chunking it into smaller pieces. Tokens can represent words, parts of words, or even individual characters depending on the \term{tokenization} method used. Most language models today use \term{Byte pair encoding} \cite{philipgageNewAlgorithmData1994} that breaks text into subwords. Almost all successful language models use probabilistic methods, therefore they are models of the form $P(t_i|t_1, t_2, ..., t_{i-1})$ where $t_i$ is the token to be predicted and $t_1, t_2, ..., t_{i-1}$ are the preceding tokens in the sequence. Earlier methods of rule based systems failed due to the complexity and ambiguity of human language and rules being too rigid \cite{NADKARNI201685,malteEvolutionTransferLearning2019}.

\subsection{Early probabilistic language models}

Early probabilistic models included simple statistical approaches such as n-grams, which relied on counting the frequency of token sequences. However this approach suffered from two competing issues, short context windows (2-grams, 3-grams etc) means that model can't understand long range dependencies in text, while larger n-grams leads to data sparsity issues where many long sequences are never seen in training data. The advent of neural networks allowed for more sophisticated models such as \term{Feedforward Neural Language models} \cite{bengioNeuralProbabilisticLanguage2003}, \term{Recurrent Neural Networks (RNNs)} \cite{tomasmikolavSTATISTICALLANGUAGEMODELS} and \term{Long Short-Term Memory networks (LSTMs)} \cite{merityRegularizingOptimizingLSTM2017} which could in theory capture longer range dependencies in text. However these models suffered from issues such as
\begin{itemize}
    \item \textbf{Vanishing gradients}: Early tokens in a long sequence have little effect on gradients and therefore the models understanding of long range dependencies.
    \item \textbf{Information bottlenecks}: The hidden state of the model has to compress all prior context into a fixed size vector.
\end{itemize}
Crucially these models failed in their flexibility to model the complex dependencies in human language.

\subsection{Modern attention based Large Language Models}

The key feature in language is context, and it is this context of surrounding words that let us humans understand and break through the complexity and ambiguity of language. The challenge is that the way in which a words company effects the meaning of the word is complex, therefore any language model needs to be able effectively capture these complex relationships. \term{Attention} is the concept where each token's representation is influenced by all the other tokens in the sequence (or just the preceding tokens in the case of \term{autoregressive} models). This was first applied to RNN-based LM in \cite{bahdanauNeuralMachineTranslation2016}, however in \cite{vaswaniAttentionAllYou2017} it was shown that attention mechanisms alone were sufficient to model language without the need for any recurrent structure\footnote{Recurrent structure is where there is a hidden state (normally a fixed length vector) that is passed along and updated each time by the next input. For language modelling this hidden state is supposed to contain the context needed to understand the next input token. However as discussed above it was not sufficient.}.

\term{Transformers} \cite{vaswaniAttentionAllYou2017} are a class of neural network \term{architecture} that utilizes the attention mechanism to process data and generate a concise representation (\term{encoder}-only), or generate data autoregressively (\term{decoder}-only). The canonical transformer architecture consists of two parts an encoder and a decoder, however decoder only architectures are shown to be sufficient for all current generative language modelling tasks. A simplified architecture of a decoder-only transformer model is shown in Figure \ref{fig:transformer_architecture}. The key feature of the transformer is the \term{self attention}. Each token has an \term{embedding vector} (used to represent its context aware meaning, i.e river \textit{bank} vs \textit{bank} robbery) which is iteratively updated by each layer of self-attention. The self-attention mechanism involves matrix multiplication between the current token and other tokens in the sequence. For generative tasks an attention mask is applied so that each token can only 'see' preceding tokens in the sequence. Furthermore to allow for different types of dependencies to be captured this attention is done multi times in parallel (\term{multi-head attention}) and merged together at the end of each layer. After several layers of self-attention the final set of token embeddings are used to generate a probability distribution over the vocabulary for the next token. A token can be sampled from this distribution, appended to the input sequence and the process repeated to generate long sequences of text.

\begin{figure}
    \centering
    \begin{tikzpicture}[
        block/.style={rectangle, draw, minimum width=3.5cm, minimum height=0.8cm, align=center, fill=white},
        arrow/.style={-{Latex[width=2mm,length=2mm]}, thick}
    ]

    % Input tokens
    \node[block] (input) {Input tokens $x_1, \dots, x_{t-1}$};

    % Embedding + Positional Encoding
    \node[block, below=0.8cm of input] (embed) {Embedding + \\ Positional Encoding};

    % --- Decoder Layer Internals ---
    % 1. Masked Self-Attention
    \node[block, below=2.4cm of embed] (attn) {Masked \\ Multi-Head Attention};

    % 2. Add & Norm (after Attention)
    \node[block, below=0.6cm of attn] (norm1) {Add \& Norm};

    % 3. Feed Forward
    \node[block, below=0.6cm of norm1] (ffn) {Feed Forward};

    % 4. Add & Norm (after FFN)
    \node[block, below=0.6cm of ffn] (norm2) {Add \& Norm};

    % Container for the repeating N layers
    \draw[dashed, thick] ([xshift=-0.6cm, yshift=0.6cm]attn.north west) rectangle ([xshift=0.6cm, yshift=-0.3cm]norm2.south east);
    \node[anchor=west] at ([xshift=0.7cm]norm1.east) {\Large $\times N$};

    % Output generation
    \node[block, below=1.5cm of norm2] (output) {Linear + Softmax \\ $\rightarrow$ Output token $x_t$};

    % --- Arrows and Connections ---
    \draw[arrow] (input) -- (embed);

    % Connection from Embed to Attention with QKV split and Residual
    % Show the standard notation: embeddings (with positional encodings) form the matrix X.
    \path (embed.south) -- (attn.north) coordinate[pos=0.35] (resid_split);
    \draw[thick] (embed.south) -- (resid_split) node[midway, right, font=\scriptsize] {$X$ is a matrix of token embeddings};
    
    % Residual 1 path (branches off before QKV split)
    \draw[thick, ->, rounded corners] (resid_split) -- ++(-3.0,0) |- (norm1.west);
    \fill (resid_split) circle (2pt);

    % QKV Split (per head)
    \coordinate (qkv_split) at ([yshift=0.4cm]attn.north);
    \draw[thick] (resid_split) -- (qkv_split) node[midway, right, align=left, font=\scriptsize] {For each head $h$: \ $Q_h = XW_Q^h$, $K_h = XW_K^h$, $V_h = XW_V^h$ \\ Attention uses $Q_h K_h^\top$ (then mask + softmax) applied to $V_h$};
    
    \draw[arrow] (qkv_split) -| ([xshift=-1.0cm]attn.north) node[pos=0.7, left, font=\scriptsize] {$V_h$};
    \draw[arrow] (qkv_split) -- (attn.north) node[pos=0.7, right, font=\scriptsize] {$K_h$};
    \draw[arrow] (qkv_split) -| ([xshift=1.0cm]attn.north) node[pos=0.7, right, font=\scriptsize] {$Q_h$};

    % Internal Layer Arrows
    \draw[arrow] (attn) -- (norm1);
    \draw[arrow] (norm1) -- (ffn);
    \draw[arrow] (ffn) -- (norm2);
    \draw[arrow] (norm2) -- (output);

    % Residual 2: Around Feed Forward
    \path (norm1.south) -- (ffn.north) coordinate[midway] (split2);
    \draw[thick, ->, rounded corners] (split2) -- ++(-3.0,0) |- (norm2.west);
    \fill (split2) circle (2pt);

    % --- Generation Loop ---
    \node[block, right=3cm of output] (gen) {Sample next token \\ and append to input};
    \draw[arrow, dashed] (output.east) -- (gen.west);
    \draw[arrow, dashed, rounded corners] (gen.north) |- (input.east);

    \end{tikzpicture}
    \caption{Simplified architecture of a decoder-only transformer model for language modeling. Adapted from \cite{vaswaniAttentionAllYou2017}.}
    \label{fig:transformer_architecture}
\end{figure}

Modern transformer based large language models (LLMs) such as GPT-3 \cite{brownLanguageModelsAre2020} are some of the largest machine learning models ever created with hundreds of billions of parameters and now trillions of parameters. Therefore training these models requires massive datasets. They are trained in a self supervised manner on large corpora (hundreds of terabytes) of text data from the Internet such as Common Crawl \cite{CommonCrawl2025}, WebText \cite{Gokaslan2019OpenWeb}, and others. The objective function is usually the cross-entropy loss between the predicted token distribution and the actual next token in the sequence. The output of this training process is a model that can generate coherent and contextually relevant text given a prompt.

\subsection{Scale of LLMs}

A LLM is a specialized type of \term{Artificial Neural Network (ANN)}. ANNs are a class of machine learning models inspired by the structure and function of biological neural networks (i.e our brain). Where a bunch of \term(neurons) are connected to gather and form a network that processes information. The key design decisions in a ANN are the architecture (which neurons are connected to which) and the \term{parameters} (the weights of the connections between neurons). As discussed above most modern LLMs use derivatives of the transformer architecture. They operate on a scale that is unprecedented in machine learning. When GPT-2 was released in 2019 it has 1.5 billion parameters \cite{radfordLanguageModelsAre2019} GPT-3 in 2020 had 175 billion parameters \cite{brownLanguageModelsAre2020} and current frontier models have in the trillions of parameters.

The scale of these models is important for two reasons. Firstly is that scale alone was found to improve the capabilities of LLMs \cite{kaplanScalingLawsNeural2020}. This is what drove the rapid increase in model size and training data over the last few years, where adding more compute was more effective than theoretical advances. Secondly is that the scale of these models means that training them requires massive computational resources. A frontier model with hundreds of billions of parameters is estimated to cost in the order of hundreds of millions to billions of dollars to train \cite{cottierRisingCostsTraining2024}. With the rise in model and dataset size there is an increasing situation where only a few very well funded organizations can train these models.

Fortunately there is an Open Source AI movement \cite{OpenSourceAI} which pushes for trained models along with the training process and data are shared publicly. This has meant that there are various \scare{open source} LLMs that reach similar capabilities to proprietary models \cite{shaoDeepSeekMathPushingLimits2024,teamGLM45AgenticReasoning2025a} \footnote{The open source movement advocates for complete transparency that includes all that would be needed to recreate the model yourself. Most of the open LLMs are simply open weight models that can have restrictive licenses.}. With these open weight models it is possible for smaller groups to fine tune them for specific tasks or domains at a fraction of the cost of training a new model from scratch. There are two key trends that can make use of these open models more effective. The first is \term{parameter efficient fine-tuning} methods \cite{huLoRALowRankAdaptation2021} that allow for effective fine-tuning of large models by only updating a small fraction of the model parameters. The second is model \term{quantization} methods \cite{dettmersQLoRAEfficientFinetuning2023,frantarGPTQAccuratePostTraining2023, xiaoSmoothQuantAccurateEfficient2024,linAWQActivationawareWeight2024,cheeQuIP2BitQuantization2024} that allow for the models parameters to be stored and computed in lower precisions formats (such as 4-bit instead of 16 or 32-bit) which drastically reduces the memory and compute requirements needed to run these models.

\subsection{LLMs in the modern world}

Large language models (LLMs) are fundamentally just language models which can predict the next token in a sequence of text. However due to the prevalence of natural language in our world and the capabilities of LLMs they have applications in a wide variety of domains. This includes Virtual Assistants, Content creation, Code generation, Translation and many more. To understand their capabilities one can look at current benchmarks and LLMs ability like winning math olympiads, passing the bar exam. The meteoric rise of capabilities has lead to problems with not having tests that are sufficiently hard \cite{phan2025humanitysexam} along with observations that on current trajectories they will outpace human thinking within the next few years \cite{danielkokotajloAI2027,measuring-ai-ability-to-complete-long-tasks}. Additionally the transformer architecture has been demonstrated to be effective in other domains like computer vision \cite{luViLBERTPretrainingTaskAgnostic2019} which has lead to \term{multi-modal models} that can process and generate text, images, and other data types \cite{alayracFlamingoVisualLanguage2022}. This flexibility of LLMs have lead to them have widespread adoption and proliferation in modern information systems \footnote{ChatGPT which was the first widely used LLM was the fastest growing website ever\cite{huChatGPTSetsRecord2023}.}. The abilities of the AI systems is such that we are seeing them employed within highly capable \term{AI agents} that perceive their environment, make plans, and act autonomously to achieve goals \cite{yaoReActSynergizingReasoning2023a,wangVoyagerOpenEndedEmbodied2023,schickToolformerLanguageModels2023,russellArtificialIntelligenceGlobal2021}.

Along with this proliferation there has been increasing concern about the safety and ethical implications of these models. Specifically these models have been shown to exhibit concerning behaviors like self-preservation and deception \cite{meinkeFrontierModelsAre2025c,greenblattAlignmentFakingLarge2024a,chenReasoningModelsDont2025a}. On top of this there are growing ethical questions around its use in our society \cite{benderDangersStochasticParrots2021, bengioManagingExtremeAI2024} as well as concerns of the welfare of the models themselves \cite{butlinConsciousnessArtificialIntelligence2023}.

\section{Challenges beyond value alignment}
\label{sec:challenges}

AI alignment is a crucial step towards safe and beneficial AI systems. There are problems that are adjacent to alignment within AI safety that need to be addressed to ensure that the introduction of AI systems into society is a net positive.

\subsection{Concrete problems of AI safety}

One can think of AI systems in an abstract way and outline the key problems that are always present when working with AI systems. There have been multiple attempts in the last decade to outline and categorize these problems \cite{amodeiConcreteProblemsAI2016a,hendrycksUnsolvedProblemsML2022}. I will outline some of the key problems here in a pseudo-categorization that can be understood as: always safe, capable and overseen.

Firstly is \textbf{robustness} that ensures the model behaves reasonably at all times, specifically in situations that are out of distribution (adversarial, black swan, shifts etc) from the training data. A key challenge to this is that models mostly are unaware of when they are out of distribution therefore they must "always behave reasonably". This is easily implemented in a vacuum robot where if it is sufficiently confused it can just stop moving, however it is unclear how to implement this level of robustness in an AI doctor or lawyer as shutting down is not a reasonable response to a difficult situation. A promising direction to solve this problem is demonstrated phenomena that  increasing model complexity increases its robustness \cite{bubeckUniversalLawRobustness2022}. At a high level this can be intuitive as a more complex model can contain and store more information about the world and therefore be able to generalize better to new situations. There has been limited evidence that this is true in modern LLM \cite{howe2025scaling}, with larger models having slightly more base robustness and being much more effective at receiving robustness training. Current models are still orders of magnitude away from the complexity of the human brain \cite{sandbergWholeBrainEmulation2008} and so may have much more "free" robustness to gain as they scale up in size and complexity.

Secondly is \textbf{factual accuracy and hallucinations} which is ensuring that the model has a correct model of the world and also a correct model of itself and its capabilities. This is particularly crucial with LLMs as they generate text which has the ability of being perfectly well written, plausible, yet completely false. This is because the model fundamentally isn't in tune with the universe and so can have no concept of when it is delusional. This is much the same problem that humans have, and we rely on external signals from society to guide our own model of how sane we are.

Lastly is \textbf{scalable oversight} which is ensuring we can keep track of and monitor the behavior of AI systems as they both scale up in capability and scale out in number of systems. If we had a provably aligned and robust system then oversight is not needed, however it is unlikely such a system will exist, instead we will likely have \scare{demonstrably} aligned and robust systems. The only way to maintain confidence that these systems are behaving well is to continually monitor them. Doing this when they are highly capable presents several challenges due to how effectively they could deceive and/or how poorly humans could even understand what they are doing. Equally so when there are a large number of systems the available resources to monitor each agent decreases and so increasingly efficient and effective monitoring techniques will be needed. There are some promising results of bootstrap supervision \cite{bowmanMeasuringProgressScalable2022,engelsScalingLawsScalable2025,kentonScalableOversightWeak2024a} where AI systems are used to help monitor other AI systems. Although potentially quite fruitful this approach introduces new challenges of misalignment propagation as one needs to ensure that the supervising AI system is also robust and aligned.

\subsection{Current negative impacts of AI}

AI is already in the world and having an effect. There are many different types of AI harms in the near term. Common themes of negative impacts of AI are as follows:
\begin{itemize}
    \item \textbf{Bias and discrimination: } AI systems can perpetuate and amplify existing biases in society. For example, in predictive policing \cite{stroudAutomatedPolicingProgram2021}, judicial sentencing \cite{juliaangwinMachineBias2016}, and hiring decisions.
    \item \textbf{Automation and job loss: } AI systems can automate tasks that were previously done by humans, leading to job loss and economic disruption \cite{autorWhyAreThere2015}. This can disproportionately affect certain groups of people, such as low-skilled workers and marginalized communities. Furthermore this leads to increased wealth concentration as the owners of companies that deploy AI systems reap the majority of the benefits.
    \item \textbf{Safety harms in high risk areas: } AI systems deployed in high-risk areas such as healthcare \cite{obermeyerDissectingRacialBias2019}, transportation, and critical infrastructure can lead to real-world safety harms if they fail or make incorrect decisions.
    \item \textbf{Population manipulation: } AI systems like recommender systems, online advertising and social media algorithms can manipulate public opinion and behavior \cite{shoshanazuboffAgeSurveillanceCapitalism2017,allcottSocialMediaFake2017}. This has had some severe real-world consequences, such as the Rohingya genocide in Myanmar \cite{UN_HRC_39_CRP2_2018}.
    \item \textbf{Explicit misuse: } More capable AI systems are at greater risk of being used for malicious purposes, which is called the dual-use problem. Common dual-use concerns include bioweapon design and cyber warfare \cite{TacklingPerilsDual2022,sandbrinkArtificialIntelligenceBiological2023}.
\end{itemize}

Many of these harms come about due to either negligent or malicious deployment of AI systems. Having systems that are better aligned to human values could help mitigate some of these harms. However there are also systemic issues with the way that AI is developed and deployed that lead to these harms. Namely the profit motive of private companies is not aligned with the goal of "helping users being their best self". Therefore solutions to these problems are the same solutions to how we stop companies in general from causing harm, namely regulation, oversight and changing the incentives of companies. This can be effectively done using the principles of public democratic AI outlined in \ref{sec:desiderata}.


\subsection{Catastrophic and existential risks of AI}

The stakes are high with AI. There are multiple plausible scenarios where AI could cause human catastrophic or existential risks to humanity. We may be living in a very vulnerable world where the majority of possible future states are negative from our current perspectives \cite{bostromVulnerableWorldHypothesis2019,manheimFragileWorldHypothesis2020}. Central to these paths and scenarios is the idea that more powerful and capable AI systems have the potential to enfeeble humans and disempower us. Once in this state of disempowerment, humans are by definition unable to control and direct their own future. Therefore, if any form of misalignment exists between the AI system and the maximum human flourishing goal, the consequences could be \term{catastrophic} or \term{existential}

AI development and AI progress has increased rapidly in the last few years \cite{kwaMeasuringAIAbility2025}. It has increased in speed due to scaling laws of AI \cite{kaplanScalingLawsNeural2020} and increased investment from both private and public sector. The key next step is that AI systems will accelerate and automate many parts of AI research. This could result in an intelligence explosion \cite{bostromSuperintelligencePathsDangers2014a,SpeculationsConcerningFirst1966} where the input that humans play in AI development tends towards zero and improvement can happen on timescales much faster than humans. Even if this intelligence explosion happens over the course of many years it is still much faster than human development over the last few thousand years. Therefore the probability of highly capable AI systems could be much higher than previously thought assuming AI can appreciably automate its own development.

With these highly capable AI systems there are several plausible catastrophic risks that could occur. These can generally be split up into two sorts of catastrophic risk, \textbf{rogue AI} and \textbf{malicious use of AI} \cite{hendrycksOverviewCatastrophicAI2023}. Rogue AI is where a system is either deceptively, overtly or indifferently doing actions against the goal of maximizing median human flourishing. The impact of this risk increases as the capabilities and proliferation of AI systems increases. The second risk also has similar increasing impact, where bad actors can do increasingly harmful things with increasingly powerful and accessible systems (bioweapons, cyber warfare etc). Either of these situations have the potential to cause an irreversible catastrophic risk to humanity. This could be through direct extinction of humanity, irreversible collapse of civilization or permanent disempowerment of humanity \cite{bostromSuperintelligencePathsDangers2014a}.

To compound these risks further are the \term{instrumental goals} that agents are likely to pursue to achieve their final goals \cite{bengioSuperintelligentAgentsPose2025a,omohundroBasicAIDrives2018}. There are three relevant instrumental goals. The first is \term{power seeking} in which any agent that wants to complete a task will be able to increase its probability of success by gaining more power and resources \cite{omohundroBasicAIDrives2018}. The second is \term{self-preservation} where being turned off will reduce the probability of success and therefore agents will try to avoid being turned off \cite{hadfield-menellOffSwitchGame2017}. The third is \term{goal-content integrity} where for an agent to achieve its goals it must ensure that its goals don't change \cite{soaresCorrigibility2015}. From these instrumental goals one could form a argument that an active intelligent system (even if perfectly aligned to maximum human flourishing) could still cause harm, much more so than an mis aligned yet much less capable system.

As the impact of these risks increases the probability of them can also be increased through both organizational incompetence (e.g. The USSR and Chernobyl, NASA and the Challenger disaster) and competitive pressures (e.g. nuclear arms race) \cite{hendrycksOverviewCatastrophicAI2023}. These two factors broadly lead to decreased safety standards with increased focus on rapid development and deployment.

\subsection{Preventing catastrophic and existential risks of AI}

Given that the stakes are very high, there should be significant effort to ensure that these risks are mitigated and negative events avoided. To prevent the risks of AI, we can either try to slow down (restrictive) or try to make it safer (constructive).

The first common approach is a halt or moratorium on AI development. Although reasonable sounding at first glance relatively large scale attempts at this have failed \cite{PauseGiantAI}. A global pause struggles from coordination problem where each actor has no incentive to pause as others might not pause and therefore gain competitive advantage. Furthermore pausing in an indefinite manner runs into the ethical problem about forgoing potential benefits of AI \footnote{This concept of whether we do need AI and the importance of bringing it to fruition compared to its potential downsides is a very important question that shouldn't be glossed over too much. However in the context of AI alignment it is outside the scope and thus not discussed further.}. This is why most pauses call for a limited time period to allow for safety, policy and governance research and implementations to catch up \cite{PauseAIProposal}. 
Given the nature of AI development requiring large amounts of resources (hardware, data, power, etc) it is plausible that a well designed pause and global coordination could be effective. It could work much like nuclear non-proliferation treaties where countries agree to limit their development of nuclear weapons and allow for inspections to ensure compliance. Also much like nuclear weapons most of current day frontier AI development requires massive data centres that are hard to hide and therefore easier to monitor, much like uranium refinement factories.

If one had time what would be the methods and ways to reduce the risks? Generally speaking they fall into technical safety solutions and governance and policy solutions. There are proposals for having various policies and regulations that would slow down the progress of AI development. Namely so that we have more time to understand the systems and develop appropriate safety measures. Policies for slowdown can involve computation restrictions, increased liability for the developers and sellers of AI systems as well as more stringent testing and evaluation requirements. Some of these can help long term safety however most of them are simply mandating a slower process and better appreciation of the risks.

With this better appreciation of the risks we can implement further policies and also develop novel technical solutions. There are various methods for AI that are safe by design \cite{dalrympleGuaranteedSafeAI2024,bengioSuperintelligentAgentsPose2025a}. The key to these methods is that the resultant AI is either deeply aligned by construction or is limited in capability to prevent misalignment from causing catastrophic risks. Another avenue towards safe super powerful AI is that they are designed to be corrigible \cite{soaresCorrigibility2015} and therefore can be corrected and modified.
